% Einheit 3 von 5 – Normale (bank/tangente-normale-schnittwinkel/e3.jsonl, zusammenbau v0.1)
%% TODO Zeitmarke Einheit 3: Marken-Zeile nicht lesbar
%% TODO Zweigzeile Teil 1: was hier gelernt wird (Satz in Schülersprache aus dem Einheitstitel) – Einheit 3
\einheitenkopf[e3]{Einheit 3 von 5 · Normale}
\zweigzeile{Abitur GK}
\setcounter{aufgabe}{17}

%% TODO Titel: Ich-kann-Satz fehlt in der Bank (Platzhalter: Kettenname) – Nr. 18
%% TODO Anweisung: ein Satz über der ersten Teilaufgabe fehlt in der Bank – Nr. 18
\begin{aufgabe}{Normale}
%% TODO \rechenplatz{n}: Schrittzahl des Grundfalls fehlt in der Bank (nur bei fester Schreibform, 2.3 b)
\begin{teile}
\teil Gesucht ist die Gerade, die den Graphen von $f$ im Punkt $P$ senkrecht schneidet; es ist $f'(x_P) = 4$. Welche Gerade, welche Steigung? \\ \kreuz{Tangente, Steigung $4$} \\ \kreuz{Normale, Steigung $-\frac{1}{4}$} \\ \kreuz{Normale, Steigung $-4$}
\teil Die Tangente an den Graphen von $f$ im Punkt $P(3 | 1)$ hat die Steigung $2$. Gleichung der Normalen in $P$? $n(x) =$ \leerfeld
\teil $f(x) = x^3 - x^2 + x + 2$; der Punkt $P(1 | y_P)$ liegt auf dem Graphen. $y_P$, Anstieg des Graphen in $P$ und Gleichung der Normalen in $P$? \hfill (FHR 2022) \\ $y_P =$ \leerfeld , Anstieg: \leerfeld , $n(x) =$ \leerfeld
\teil Die Tangente $t$ hat die Gleichung $y = 2x + 1$. Gleichung der Geraden durch $Q(4 | 1)$, die senkrecht zu $t$ verläuft? $y =$ \leerfeld
\teil Ein Bogen wird durch $f(x) = \frac{1}{8}x^2 + 1$ beschrieben ($1$ LE = $1$ m, die $x$-Achse ist der Boden). Eine Strebe führt vom Punkt $P(4 | f(4))$ senkrecht zur Tangente in $P$ bis zum Boden. Wie lang ist sie? \hfill (Abitur 2024 LK) \\ \leerfeld[m]
\teil Abstand des Punktes $Q(10 | 0)$ von der Geraden $y = 3x$ – Lotfußpunkt und Abstand? \leerfeld
\teil Für eine Funktion $f$ gilt $f'(x) = \frac{1}{2}x^2$. Ein Kreis mit Mittelpunkt auf der $y$-Achse berührt den Graphen von $f'$ im Punkt $A(2 | 2)$. $y$-Koordinate des Mittelpunkts, rechnerisch? \hfill (Abitur 2026 LK) \\ $y_M =$ \leerfeld
\teil Für ein Hangprofil $h$ wird die Gleichung $\frac{h(x) - 2}{x - 3} \cdot h'(x) = -1$ gelöst. Welcher Abstand lässt sich aus der Lösung berechnen? Begründe. \hfill (Abitur 2023 LK)
\end{teile}
\end{aufgabe}

%% TODO Titel: Ich-kann-Satz fehlt in der Bank (Platzhalter: Kettenname) – Nr. 19
%% TODO Anweisung: ein Satz über der ersten Teilaufgabe fehlt in der Bank – Nr. 19
\begin{aufgabe}{Normale – Fehler finden, Begründen}
\begin{teile}
\teil Ich finde den Fehler: Die Tangente im Punkt $P(3 | 2)$ hat die Steigung $3$. Mira bestimmt die Normale: \rechnung{m_n &= -3 \\ n(x) &= -3x + 11}
\teil Begründe, warum der kürzeste Abstand eines Punktes von einer Geraden auf dem Lot liegt.
\end{teile}
\end{aufgabe}
